\section{Sim-to-Real y experimentos de verificación}

\subsection{Calibración de la simulación y verificación en el mundo real}

\begin{figure}[H]
\centering
\includegraphics[width=0.8\textwidth]{slide-02-11.jpg}
\caption{Slide 02-11: cómo se calibra un pipeline de sim-to-real mediante domain randomization y réplicas del mundo real.}
\end{figure}

Para que la policy migre de la simulación al entorno real, las estrategias habituales incluyen domain randomization, dynamics randomization y fidelity tuning. Cada migración debe registrar los parámetros de randomization y las métricas de monitoreo del reality gap.

\begin{table}[H]
\centering
\begin{tabular}{p{0.3\textwidth} p{0.6\textwidth}}
\toprule
Punto de verificación & Descripción \\
\midrule
Domain gap & distribution drift score de las diferencias de dinámica/percepción \\
Simulator fidelity & error L2 entre la fidelity visual/de control y el real world log \\
Transfer success & task success rate del rollout en el mundo real \\
\bottomrule
\end{tabular}
\caption{Matriz de verificación de Sim-to-Real}
\end{table}

\begin{knowledgebox}{Registrar los metadatos de transfer}
Usar metadata para registrar, en cada transfer, el simulator seed, el randomization range y la hardware configuration ayuda a reproducir el éxito o el failure de un salto determinado.
\end{knowledgebox}

\subsection{Cross-environment evaluation}

Distintos entornos (lab, field, cloud) tienen distintos constraints. Es necesario fijar un evaluation order, por ejemplo: primero ejecutar el benchmark en el lab environment, después revisar el sensor noise y el drift en el field environment, y por último recolectar logs en el cloud pipeline.

\begin{warningbox}{No usar el baseline directamente en field rollouts}
Una vez que el field environment sufre drift, resulta muy expensive. Hay que verificar primero la métrica en el lab, luego hacer una incremental expansion hacia el field, y por último recolectar los logs en el cloud dentro de un evidence board.
\end{warningbox}

\subsection{Resumen de la sección}

El pipeline de Sim-to-Real depende de una structured validation matrix y de metadata para poder migrar la policy con seguridad entre distintos entornos; antes de cada field rollout es indispensable verificar el domain gap.
